% preamble-pl.tex -- sdílená preamble pro pracovní listy
% Includuje se z ucitel/lekceXX/pracovni_list.tex jako:
%   % preamble-pl.tex -- sdílená preamble pro pracovní listy
% Includuje se z ucitel/lekceXX/pracovni_list.tex jako:
%   % preamble-pl.tex -- sdílená preamble pro pracovní listy
% Includuje se z ucitel/lekceXX/pracovni_list.tex jako:
%   % preamble-pl.tex -- sdílená preamble pro pracovní listy
% Includuje se z ucitel/lekceXX/pracovni_list.tex jako:
%   \input{../spolecne/preamble-pl.tex}

\documentclass[a4paper,12pt]{article}

% --- Jazyk a font ---
\usepackage{polyglossia}
\setmainlanguage{czech}
\usepackage{fontspec}

% --- Okraje ---
\usepackage[top=20mm, bottom=20mm, left=22mm, right=22mm]{geometry}

% --- Česky korektní uvozovky ---
\usepackage{csquotes}
\newcommand{\uv}[1]{\enquote{#1}}

% --- Typografie ---
\usepackage{microtype}
\usepackage{parskip}          % odstavce odděleny mezerou, ne odsazením

% --- Barvy a grafika ---
\usepackage[table]{xcolor}
\definecolor{microbit-blue}{HTML}{1E4D8C}
\definecolor{code-bg}{HTML}{F5F5F5}
\definecolor{task-bg}{HTML}{EEF4FF}

% --- Tabulky ---
\usepackage{booktabs}
\usepackage{multirow}

% --- Zdrojový kód (Python) ---
\usepackage{listings}
\lstset{
  language=Python,
  basicstyle=\ttfamily\small,
  backgroundcolor=\color{code-bg},
  frame=single,
  framerule=0pt,
  rulecolor=\color{gray!30},
  xleftmargin=6pt,
  xrightmargin=6pt,
  breaklines=true,
  showstringspaces=false,
  keywordstyle=\color{microbit-blue}\bfseries,
  stringstyle=\color{teal},
  commentstyle=\color{gray}\itshape,
  literate=
    {á}{{\'a}}1 {é}{{\'e}}1 {í}{{\'i}}1 {ó}{{\'o}}1 {ú}{{\'u}}1
    {ů}{{\r{u}}}1 {č}{{\v{c}}}1 {š}{{\v{s}}}1 {ž}{{\v{z}}}1
    {Á}{{\'A}}1 {É}{{\'E}}1 {Í}{{\'I}}1 {Ó}{{\'O}}1 {Ú}{{\'U}}1
    {Č}{{\v{C}}}1 {Š}{{\v{S}}}1 {Ž}{{\v{Z}}}1 {ň}{{\v{n}}}1 {ř}{{\v{r}}}1,
}

% --- Zadání úkolů ---
\usepackage[most]{tcolorbox}
\tcbuselibrary{breakable}

\newcounter{ukol}[section]
\newenvironment{ukol}[1][]{%
  \stepcounter{ukol}%
  \begin{tcolorbox}[
    colback=task-bg,
    colframe=microbit-blue,
    fonttitle=\bfseries,
    title={Úkol~\theukol\ifx\relax#1\relax\else{: #1}\fi},
    breakable,
    left=6pt, right=6pt, top=4pt, bottom=4pt
  ]%
}{%
  \end{tcolorbox}%
}

% Inline kód — underscore package zajistí, ze _ funguje v texttt
\usepackage{underscore}
\newcommand{\pyinline}[1]{\texttt{#1}}
% Pojem (zvýraznění termínu) — stejné jako v preamble-prez.tex
\newcommand{\pojem}[1]{\textbf{#1}}

% Místo pro odpověď studenta
\newcommand{\odpoved}[1][3cm]{%
  \vspace{#1}%
}

% --- Záhlaví / zápatí ---
\usepackage{fancyhdr}
\pagestyle{fancy}
\fancyhf{}
\renewcommand{\headrulewidth}{0.4pt}
\fancyhead[L]{\small\textcolor{microbit-blue}{\textbf{Úvod do Pythonu na Micro:bitu}}}
\fancyhead[R]{\small\textcolor{gray}{\lekce}}
\fancyfoot[C]{\small\thepage}

% --- Ostatní ---
\usepackage{enumitem}
\usepackage{graphicx}
\usepackage{hyperref}
\hypersetup{
  colorlinks=true,
  linkcolor=microbit-blue,
  urlcolor=microbit-blue,
}


\documentclass[a4paper,12pt]{article}

% --- Jazyk a font ---
\usepackage{polyglossia}
\setmainlanguage{czech}
\usepackage{fontspec}

% --- Okraje ---
\usepackage[top=20mm, bottom=20mm, left=22mm, right=22mm]{geometry}

% --- Česky korektní uvozovky ---
\usepackage{csquotes}
\newcommand{\uv}[1]{\enquote{#1}}

% --- Typografie ---
\usepackage{microtype}
\usepackage{parskip}          % odstavce odděleny mezerou, ne odsazením

% --- Barvy a grafika ---
\usepackage[table]{xcolor}
\definecolor{microbit-blue}{HTML}{1E4D8C}
\definecolor{code-bg}{HTML}{F5F5F5}
\definecolor{task-bg}{HTML}{EEF4FF}

% --- Tabulky ---
\usepackage{booktabs}
\usepackage{multirow}

% --- Zdrojový kód (Python) ---
\usepackage{listings}
\lstset{
  language=Python,
  basicstyle=\ttfamily\small,
  backgroundcolor=\color{code-bg},
  frame=single,
  framerule=0pt,
  rulecolor=\color{gray!30},
  xleftmargin=6pt,
  xrightmargin=6pt,
  breaklines=true,
  showstringspaces=false,
  keywordstyle=\color{microbit-blue}\bfseries,
  stringstyle=\color{teal},
  commentstyle=\color{gray}\itshape,
  literate=
    {á}{{\'a}}1 {é}{{\'e}}1 {í}{{\'i}}1 {ó}{{\'o}}1 {ú}{{\'u}}1
    {ů}{{\r{u}}}1 {č}{{\v{c}}}1 {š}{{\v{s}}}1 {ž}{{\v{z}}}1
    {Á}{{\'A}}1 {É}{{\'E}}1 {Í}{{\'I}}1 {Ó}{{\'O}}1 {Ú}{{\'U}}1
    {Č}{{\v{C}}}1 {Š}{{\v{S}}}1 {Ž}{{\v{Z}}}1 {ň}{{\v{n}}}1 {ř}{{\v{r}}}1,
}

% --- Zadání úkolů ---
\usepackage[most]{tcolorbox}
\tcbuselibrary{breakable}

\newcounter{ukol}[section]
\newenvironment{ukol}[1][]{%
  \stepcounter{ukol}%
  \begin{tcolorbox}[
    colback=task-bg,
    colframe=microbit-blue,
    fonttitle=\bfseries,
    title={Úkol~\theukol\ifx\relax#1\relax\else{: #1}\fi},
    breakable,
    left=6pt, right=6pt, top=4pt, bottom=4pt
  ]%
}{%
  \end{tcolorbox}%
}

% Inline kód — underscore package zajistí, ze _ funguje v texttt
\usepackage{underscore}
\newcommand{\pyinline}[1]{\texttt{#1}}
% Pojem (zvýraznění termínu) — stejné jako v preamble-prez.tex
\newcommand{\pojem}[1]{\textbf{#1}}

% Místo pro odpověď studenta
\newcommand{\odpoved}[1][3cm]{%
  \vspace{#1}%
}

% --- Záhlaví / zápatí ---
\usepackage{fancyhdr}
\pagestyle{fancy}
\fancyhf{}
\renewcommand{\headrulewidth}{0.4pt}
\fancyhead[L]{\small\textcolor{microbit-blue}{\textbf{Úvod do Pythonu na Micro:bitu}}}
\fancyhead[R]{\small\textcolor{gray}{\lekce}}
\fancyfoot[C]{\small\thepage}

% --- Ostatní ---
\usepackage{enumitem}
\usepackage{graphicx}
\usepackage{hyperref}
\hypersetup{
  colorlinks=true,
  linkcolor=microbit-blue,
  urlcolor=microbit-blue,
}


\documentclass[a4paper,12pt]{article}

% --- Jazyk a font ---
\usepackage{polyglossia}
\setmainlanguage{czech}
\usepackage{fontspec}

% --- Okraje ---
\usepackage[top=20mm, bottom=20mm, left=22mm, right=22mm]{geometry}

% --- Česky korektní uvozovky ---
\usepackage{csquotes}
\newcommand{\uv}[1]{\enquote{#1}}

% --- Typografie ---
\usepackage{microtype}
\usepackage{parskip}          % odstavce odděleny mezerou, ne odsazením

% --- Barvy a grafika ---
\usepackage[table]{xcolor}
\definecolor{microbit-blue}{HTML}{1E4D8C}
\definecolor{code-bg}{HTML}{F5F5F5}
\definecolor{task-bg}{HTML}{EEF4FF}

% --- Tabulky ---
\usepackage{booktabs}
\usepackage{multirow}

% --- Zdrojový kód (Python) ---
\usepackage{listings}
\lstset{
  language=Python,
  basicstyle=\ttfamily\small,
  backgroundcolor=\color{code-bg},
  frame=single,
  framerule=0pt,
  rulecolor=\color{gray!30},
  xleftmargin=6pt,
  xrightmargin=6pt,
  breaklines=true,
  showstringspaces=false,
  keywordstyle=\color{microbit-blue}\bfseries,
  stringstyle=\color{teal},
  commentstyle=\color{gray}\itshape,
  literate=
    {á}{{\'a}}1 {é}{{\'e}}1 {í}{{\'i}}1 {ó}{{\'o}}1 {ú}{{\'u}}1
    {ů}{{\r{u}}}1 {č}{{\v{c}}}1 {š}{{\v{s}}}1 {ž}{{\v{z}}}1
    {Á}{{\'A}}1 {É}{{\'E}}1 {Í}{{\'I}}1 {Ó}{{\'O}}1 {Ú}{{\'U}}1
    {Č}{{\v{C}}}1 {Š}{{\v{S}}}1 {Ž}{{\v{Z}}}1 {ň}{{\v{n}}}1 {ř}{{\v{r}}}1,
}

% --- Zadání úkolů ---
\usepackage[most]{tcolorbox}
\tcbuselibrary{breakable}

\newcounter{ukol}[section]
\newenvironment{ukol}[1][]{%
  \stepcounter{ukol}%
  \begin{tcolorbox}[
    colback=task-bg,
    colframe=microbit-blue,
    fonttitle=\bfseries,
    title={Úkol~\theukol\ifx\relax#1\relax\else{: #1}\fi},
    breakable,
    left=6pt, right=6pt, top=4pt, bottom=4pt
  ]%
}{%
  \end{tcolorbox}%
}

% Inline kód — underscore package zajistí, ze _ funguje v texttt
\usepackage{underscore}
\newcommand{\pyinline}[1]{\texttt{#1}}
% Pojem (zvýraznění termínu) — stejné jako v preamble-prez.tex
\newcommand{\pojem}[1]{\textbf{#1}}

% Místo pro odpověď studenta
\newcommand{\odpoved}[1][3cm]{%
  \vspace{#1}%
}

% --- Záhlaví / zápatí ---
\usepackage{fancyhdr}
\pagestyle{fancy}
\fancyhf{}
\renewcommand{\headrulewidth}{0.4pt}
\fancyhead[L]{\small\textcolor{microbit-blue}{\textbf{Úvod do Pythonu na Micro:bitu}}}
\fancyhead[R]{\small\textcolor{gray}{\lekce}}
\fancyfoot[C]{\small\thepage}

% --- Ostatní ---
\usepackage{enumitem}
\usepackage{graphicx}
\usepackage{hyperref}
\hypersetup{
  colorlinks=true,
  linkcolor=microbit-blue,
  urlcolor=microbit-blue,
}


\documentclass[a4paper,12pt]{article}

% --- Jazyk a font ---
\usepackage{polyglossia}
\setmainlanguage{czech}
\usepackage{fontspec}

% --- Okraje ---
\usepackage[top=20mm, bottom=20mm, left=22mm, right=22mm]{geometry}

% --- Česky korektní uvozovky ---
\usepackage{csquotes}
\newcommand{\uv}[1]{\enquote{#1}}

% --- Typografie ---
\usepackage{microtype}
\usepackage{parskip}          % odstavce odděleny mezerou, ne odsazením

% --- Barvy a grafika ---
\usepackage[table]{xcolor}
\definecolor{microbit-blue}{HTML}{1E4D8C}
\definecolor{code-bg}{HTML}{F5F5F5}
\definecolor{task-bg}{HTML}{EEF4FF}

% --- Tabulky ---
\usepackage{booktabs}
\usepackage{multirow}

% --- Zdrojový kód (Python) ---
\usepackage{listings}
\lstset{
  language=Python,
  basicstyle=\ttfamily\small,
  backgroundcolor=\color{code-bg},
  frame=single,
  framerule=0pt,
  rulecolor=\color{gray!30},
  xleftmargin=6pt,
  xrightmargin=6pt,
  breaklines=true,
  showstringspaces=false,
  keywordstyle=\color{microbit-blue}\bfseries,
  stringstyle=\color{teal},
  commentstyle=\color{gray}\itshape,
  literate=
    {á}{{\'a}}1 {é}{{\'e}}1 {í}{{\'i}}1 {ó}{{\'o}}1 {ú}{{\'u}}1
    {ů}{{\r{u}}}1 {č}{{\v{c}}}1 {š}{{\v{s}}}1 {ž}{{\v{z}}}1
    {Á}{{\'A}}1 {É}{{\'E}}1 {Í}{{\'I}}1 {Ó}{{\'O}}1 {Ú}{{\'U}}1
    {Č}{{\v{C}}}1 {Š}{{\v{S}}}1 {Ž}{{\v{Z}}}1 {ň}{{\v{n}}}1 {ř}{{\v{r}}}1,
}

% --- Zadání úkolů ---
\usepackage[most]{tcolorbox}
\tcbuselibrary{breakable}

\newcounter{ukol}[section]
\newenvironment{ukol}[1][]{%
  \stepcounter{ukol}%
  \begin{tcolorbox}[
    colback=task-bg,
    colframe=microbit-blue,
    fonttitle=\bfseries,
    title={Úkol~\theukol\ifx\relax#1\relax\else{: #1}\fi},
    breakable,
    left=6pt, right=6pt, top=4pt, bottom=4pt
  ]%
}{%
  \end{tcolorbox}%
}

% Inline kód — underscore package zajistí, ze _ funguje v texttt
\usepackage{underscore}
\newcommand{\pyinline}[1]{\texttt{#1}}
% Pojem (zvýraznění termínu) — stejné jako v preamble-prez.tex
\newcommand{\pojem}[1]{\textbf{#1}}

% Místo pro odpověď studenta
\newcommand{\odpoved}[1][3cm]{%
  \vspace{#1}%
}

% --- Záhlaví / zápatí ---
\usepackage{fancyhdr}
\pagestyle{fancy}
\fancyhf{}
\renewcommand{\headrulewidth}{0.4pt}
\fancyhead[L]{\small\textcolor{microbit-blue}{\textbf{Úvod do Pythonu na Micro:bitu}}}
\fancyhead[R]{\small\textcolor{gray}{\lekce}}
\fancyfoot[C]{\small\thepage}

% --- Ostatní ---
\usepackage{enumitem}
\usepackage{graphicx}
\usepackage{hyperref}
\hypersetup{
  colorlinks=true,
  linkcolor=microbit-blue,
  urlcolor=microbit-blue,
}
